\begin{figure}[t]
    \centering
    \includegraphics[width=1.\linewidth, trim={0 0 0 0}, clip]{figure/pdf/rays.pdf}
    \vspace{-1.5em}
    \caption{
        \textbf{Comparison of camera encoding methods.}
        (a)~\emph{Direct Parameterization} encodes camera intrinsics and extrinsics as raw parameters, which lacks geometric interpretability and compatibility across camera types.
        (b)~\emph{Pl\"ucker Encoding} represents each ray as a pair of direction and moment vectors, providing a physically grounded but absolute, coordinate-dependent description.
        (c)~\emph{Projective Positional Encoding} encodes relative cameras in projective space, yet assumes pinhole projection and cannot model non-linear lens distortions.
        (d)~Our \emph{Relative Ray Encoding} reformulates geometric relationships in ray space, where each token corresponds to its own viewing ray, enabling better pose generalization and compatibility with \textbf{arbitrary} camera lenses.
    }\label{fig:rays}
    \vspace{-1em}
\end{figure}
